\documentclass[12pt]{article}
% -------------------
% Packages
% -------------------
\usepackage{
	amsmath,			% Math Environments
	amssymb,			% Extended Symbols
	amsthm,			    % Theorem Environments
	cancel,			    % Use Cancels
	enumerate,		    % Enumerate Environments
	graphicx,			% Include Images
	lastpage,			% Reference Lastpage
	multicol,			% Use Multi-columns
	multirow,			% Use Multi-rows
	xcolor,			    % Use Colors
	float,
	geometry,
	quiver
	}

\usetikzlibrary{nfold}

% -------------------
% Header & Footer
% -------------------
\usepackage{fancyhdr}

\newcommand{\assignment}[1]{
	\fancypagestyle{title}{
		%Headers
		\fancyhead[L]{Jake Bridge}
		\fancyhead[C]{MATH 418}
		\fancyhead[R]{HW #1}
		\renewcommand{\headrulewidth}{0.2pt}
		%Footers
		\fancyfoot[L]{}
		\fancyfoot[C]{}
		\fancyfoot[R]{}
		\renewcommand{\footrulewidth}{0.0pt}
	}
	\thispagestyle{title}
}
\fancypagestyle{pages}{
	%Headers
	\fancyhead[L]{}
	\fancyhead[C]{}
	\fancyhead[R]{}
	\renewcommand{\headrulewidth}{0.0pt}
	%Footers
	\fancyfoot[L]{}
	\fancyfoot[C]{}
	\fancyfoot[R]{}
	\renewcommand{\footrulewidth}{0.0pt}
}
\headheight=18pt
\footskip=14pt

\pagestyle{pages}

% ---------------
% Page Layout
% ---------------
\geometry{letterpaper,
bindingoffset=0.2in,
left=1in,
right=1in,
top=1in,
bottom=1in,
footskip=.25in}
% -------------------
% Font
% -------------------
%\usepackage[T1]{fontenc}
%\usepackage{charter}

%\usepackage[T1]{fontenc}
%\usepackage{mathpazo}

%\usepackage[bitstream-charter]{mathdesign}
%\usepackage[T1]{fontenc}


% -------------------
% Commands
% -------------------

% Problem Labels
\newcounter{problem}[section]
\newenvironment{problem}[1][\theproblem]{\refstepcounter{problem}\noindent \textbf{Problem~#1.\quad}}{\vskip\smallskipamount}


\newenvironment{solution}{\noindent \textbf{Solution.\quad}}{\vskip\bigskipamount}


% Special Characters
\newcommand{\N}{\mathbb{N}}
\newcommand{\Z}{\mathbb{Z}}
\newcommand{\Q}{\mathbb{Q}}
\newcommand{\R}{\mathbb{R}}
\newcommand{\C}{\mathbb{C}}

\newcommand{\mc}[1]{\mathcal{#1}}

\newcommand{\ob}[1]{\text{ob}(#1)}
\newcommand{\natt}{\Rightarrow}
\newcommand{\iso}{\cong}
\newcommand{\op}{\text{op}}

\newcommand{\Set}{\textbf{Set}}
\newcommand{\Hom}{\textbf{Hom}}
\newcommand{\Small}{\mathbb{S}}

\newcommand{\eps}{\varepsilon}

% Math Operators

% Special Commands



% DOC
\begin{document}
	\assignment{9}

\begin{problem}[87]
	Prove that any equivalence of categories preserves all limits (that exist in the source)
\end{problem}
\begin{solution}
	Let $F: \mc{A} \to \mc{B}$ be an equivalence of categories (in the sense that $F$ is an essentially surjective fully faithful functor). Let $\Small$ be a small category and let $D: \Small \to \mc{A}$ be a diagram. Suppose $(A, f_S: A \to DS)$ is a limit cone of $D$ in $\mc{A}$. 
	
	We must show that $(FA, Ff_S : FA \to FDS)$ is a limit cone in $\mc{B}$. It is clearly a cone by functoriality of $F$. Let $(B, g_S: B \to FDS)$ be another cone. $F$ is full, faithful, and essentially surjective. Essential surjectivity tells us that $B$ is isomorphic to some $FA', A' \in \mc{C}$. Call the isomorphism $\alpha: FA' \to B$. Hence, we have a morphism $g_S \circ \alpha: FA' \to FDS$. By full and faithfulness of $F$, we get a corresponding morphism $\widehat{g_S \circ \alpha}: A' \to DS$. But, by $A$ being a cone, we find a unique $\hat{h}: A' \to A$ such that $\widehat{g_S \circ \alpha}$ factors through $\hat{h}$. We can push this through $F$ to get $h: FA' \to FA$ such that $g_S \circ \alpha = Ff_S \circ h$. 
	
	% https://q.uiver.app/#q=WzAsNCxbMiwxLCJGQSJdLFswLDEsIkIiXSxbMiwyLCJGRFMiXSxbMCwwLCJGQSciXSxbMSwwLCJoIFxcY2lyYyBcXGFscGhhXnstMX0iLDEseyJzdHlsZSI6eyJib2R5Ijp7Im5hbWUiOiJkYXNoZWQifX19XSxbMCwyLCJGZl9TIl0sWzEsMiwiZ19TIiwxXSxbMywxLCJcXGFscGhhIFxcaXNvICIsMl0sWzMsMCwiXFxleGlzdHMhIGgiLDEseyJzdHlsZSI6eyJib2R5Ijp7Im5hbWUiOiJkYXNoZWQifX19XV0=
	\[\begin{tikzcd}
		{FA'} && \\
		B && FA \\
		&& FDS
		\arrow["{\alpha \iso }"', from=1-1, to=2-1]
		\arrow["{\exists! h}"{description}, dashed, from=1-1, to=2-3]
		\arrow["{h \circ \alpha^{-1}}"{description}, dashed, from=2-1, to=2-3]
		\arrow["{g_S}"{description}, from=2-1, to=3-3]
		\arrow["{Ff_S}", from=2-3, to=3-3]
	\end{tikzcd}\]
	
	Then, we get that $g_S$ factors uniquely as $Ff_S \circ (h \circ \alpha^{-1})$. Hence, $(FA, Ff_S)$ is a limit cone in $\mc{B}$. So, $F$ preserves limits! \qed
\end{solution}





\begin{problem}[91]
	Let $\mc{C}$ be a category with two objects and all products. Does it have an initial object? A terminal object?
\end{problem}
\begin{solution}
	
	Let $\Small$ be the empty category. Then, the empty product is the limit over $D: \Small \to \mc{C}$, which is a terminal object. 
	
	With regard to initial objects, I don't have an answer, but will give some thoughts...
	
	If we have a unique isomorphism between $A$ and $B$, and $B$ is terminal, then $A$ is terminal, so they would both also be initial. That's the easy way out. I've not been able to prove this, and don't think its true. The other way would be to have some $\Hom(A, A)$ with a ton of morphisms and a unique morphism $A \to B$ (WLOG $B$ is the terminal object). To show that $A$ is initial, we'd need to show that the only morphism in that hom set is the identity. We can do that pretty easily if we have only isomorphisms (I take an even simpler example below).
	
	Take, for example, $A = \{1, 2\}, B = \{0\}$ with morphisms $g: A \to B$ that takes $g(1) = 0, g(2) = 0$ and $s: A \to A$ that takes $s(1) = 2, s(2) = 1$ (and the identities, of course). $B$ is a terminal element, (hence empty product), and $A$ and $B$ themselves are single products. Any product containing a copy of $A$ must be $A$, as otherwise, we don't get any projection maps $B\to A$, which are required for a product. We know $s \circ g = 1 \circ g$.

	
	This tells us, though, that there is no choice of $\pi_1$ and $\pi_2$ to make the binary product work. Indeed, if both are $1$, the choices $s, 1$ don't give a unique morphism. If they are both $s$, the same. If they are different (WLOG $\pi_1 = s, \pi_2 = 1$, then the choices $s, s$ or $1, 1$ don't yield a unique morphism). 
	
	I'd like to use this as motivation...
	
	Suppose that our category is locally small and admits all products. WLOG, as above, $B$ is terminal. Then, given any set $I$ of morphisms $A \to A$, we construct the product of $A$ over $I$, which must be $A$ to have projection maps, as above. (My notation will suggest that $I$ is countable. This is not necessarily the case). Then, we get projection maps $\pi_i$. If two projection maps are $1$, choose different maps for those two. This will break the product structure. So, we can have at most one map that is the identity. 
	
	Without supposing isomorphisms, though, I'm not sure I have an easy way to guarantee we can build some cone that breaks the product, which leads me to believe that we are not guaranteed an initial object. I don't have an example, yet, though---in $\Set$, at least, to have all products between only two objects we need both to be one element sets, and we get to the boring case where the objects are both initial and terminal in the full subcategory of $\Set$ with the two singletons. 
	
	My suspicion is that for a locally small 2 object category, we need some strange endomorphisms to get every product and not have an initial object, but it could well be possible. 
	
	Or, perhaps there's a very simple example I've overlooked. Who knows? Not me!
	
	
	
	
	
	
	There is another interesting question, as well, if we suppose we only have binary products, that I will not answer here...
	
%	% https://q.uiver.app/#q=WzAsNCxbMSwyLCJBIl0sWzIsMywiQiJdLFswLDMsIkEiXSxbMSwwLCJBIl0sWzAsMiwiXFxwaV9BIiwxXSxbMCwxLCJnIiwxXSxbMywyLCIxX0EiLDJdLFszLDAsIlxcZXhpc3RzISIsMSx7InN0eWxlIjp7ImJvZHkiOnsibmFtZSI6ImRhc2hlZCJ9fX1dLFszLDEsImciLDFdXQ==
%	\[\begin{tikzcd}
%		& A & \\
%		\\
%		& A \\
%		A && B
%		\arrow["{\exists!}"{description}, dashed, from=1-2, to=3-2]
%		\arrow["{1_A}"', from=1-2, to=4-1]
%		\arrow["g"{description}, from=1-2, to=4-3]
%		\arrow["{\pi_A}"{description}, from=3-2, to=4-1]
%		\arrow["g"{description}, from=3-2, to=4-3]
%	\end{tikzcd}\]
%	
%	In general, this also tells us that $
%	
%	Then, this tells us that there is a unique $f_1$ such that $\pi_A = \pi_A \circ f_1$. Taking inverses, we get that $f_1 = 1_A$. 
%	% https://q.uiver.app/#q=WzAsNCxbMSwyLCJBIl0sWzIsMywiQiJdLFswLDMsIkEiXSxbMSwwLCJBIl0sWzAsMiwiXFxwaV9BIiwxXSxbMCwxLCJnIiwxXSxbMywyLCJcXHBpX0EiLDJdLFszLDAsIlxcZXhpc3RzISIsMSx7InN0eWxlIjp7ImJvZHkiOnsibmFtZSI6ImRhc2hlZCJ9fX1dLFszLDEsImciLDFdXQ==
%	\[\begin{tikzcd}
%		& A & \\
%		\\
%		& A \\
%		A && B
%		\arrow["{\exists!}"{description}, dashed, from=1-2, to=3-2]
%		\arrow["{\pi_A}"', from=1-2, to=4-1]
%		\arrow["g"{description}, from=1-2, to=4-3]
%		\arrow["{\pi_A}"{description}, from=3-2, to=4-1]
%		\arrow["g"{description}, from=3-2, to=4-3]
%	\end{tikzcd}\]
	
	

	
	
	
%	
%	Then, if $A \times A = A$, we know there is a unique map $g$ such that $\pi_1 \circ g = 1, \pi_2 \circ g = 1$. 
%	
%	If $\mc{C}$ only has binary products, we have that $A \times B = A$ (WLOG).
%	
%	We will show that $A$ is initial: that there is a unique morphism from $A \to A$ and another unique morphism from $A \to B$. These morphisms will be the identity and projection map, respectively. 
%	
%	% https://q.uiver.app/#q=WzAsNCxbMSwwLCJBIl0sWzEsMiwiQSJdLFswLDMsIkEiXSxbMiwzLCJCIl0sWzAsMSwiMV9BIiwxLHsic3R5bGUiOnsiYm9keSI6eyJuYW1lIjoiZGFzaGVkIn19fV0sWzAsMywiXFxwaV8yIl0sWzEsMiwiXFxwaV8xIl0sWzEsMywiXFxwaV8yIiwyXSxbMCwyLCIxX0EiLDJdXQ==
%	\[\begin{tikzcd}
%		& A & \\
%		\\
%		& A \\
%		A && B
%		\arrow["{1_A}"{description}, dashed, from=1-2, to=3-2]
%		\arrow["{1_A}"', from=1-2, to=4-1]
%		\arrow["{\pi_2}", from=1-2, to=4-3]
%		\arrow["{\pi_1}", from=3-2, to=4-1]
%		\arrow["{\pi_2}"', from=3-2, to=4-3]
%	\end{tikzcd}\]
%	
%	This diagram tells us that $\pi_1 = 1_A$. We need to show that $\pi_2$ is the only morphism from $A \to B$ and that there cannot be any more morphisms $A\to A$. Let $f: A \to B$ be another map. Uniqueness of the identity gives that % https://q.uiver.app/#q=WzAsNCxbMSwwLCJBIl0sWzEsMiwiQSJdLFswLDMsIkEiXSxbMiwzLCJCIl0sWzAsMSwiMV9BIiwxLHsic3R5bGUiOnsiYm9keSI6eyJuYW1lIjoiZGFzaGVkIn19fV0sWzAsMywiZiJdLFsxLDIsIjFfQSJdLFsxLDMsIlxccGlfMiIsMl0sWzAsMiwiMV9BIiwyXV0=
%	\[\begin{tikzcd}
%		& A & \\
%		\\
%		& A \\
%		A && B
%		\arrow["{1_A}"{description}, dashed, from=1-2, to=3-2]
%		\arrow["{1_A}"', from=1-2, to=4-1]
%		\arrow["f", from=1-2, to=4-3]
%		\arrow["{1_A}", from=3-2, to=4-1]
%		\arrow["{\pi_2}"', from=3-2, to=4-3]
%	\end{tikzcd}\]
%	commutes. Hence, $\pi_2$ is unique. 
%	
%	Lastly, if $g: A \to A$ is a morphism, % https://q.uiver.app/#q=WzAsNCxbMSwwLCJBIl0sWzEsMiwiQSJdLFswLDMsIkEiXSxbMiwzLCJCIl0sWzAsMSwiMV9BIiwxLHsic3R5bGUiOnsiYm9keSI6eyJuYW1lIjoiZGFzaGVkIn19fV0sWzAsMywiXFxwaV8yIl0sWzEsMiwiMV9BIl0sWzEsMywiXFxwaV8yIiwyXSxbMCwyLCJnIiwyXV0=
%	\[\begin{tikzcd}
%		& A & \\
%		\\
%		& A \\
%		A && B
%		\arrow["{g}"{description}, dashed, from=1-2, to=3-2]
%		\arrow["g"', from=1-2, to=4-1]
%		\arrow["{\pi_2}", from=1-2, to=4-3]
%		\arrow["{1_A}", from=3-2, to=4-1]
%		\arrow["{\pi_2}"', from=3-2, to=4-3]
%	\end{tikzcd}\]
%	tells us that $\pi_2 \circ g = \pi_2$. 	By uniqueness of morphisms, $g = 1_A$, so there is a unique morphism $1_A: A \to A$ and unique morphism $\pi_2: A\to B$. This is the same as saying that $A \times B = A$ is initial. 
%	
%	
%	
% We know that $\pi_2 \circ g = 1_B$ by the universal property of the product again: 
%	% https://q.uiver.app/#q=WzAsNCxbMSwwLCJCIl0sWzEsMSwiQSJdLFswLDIsIkEiXSxbMiwyLCJCIl0sWzEsMiwiMV9BIl0sWzEsMywiXFxwaV8yIiwyXSxbMCwxLCJnIiwyLHsic3R5bGUiOnsiYm9keSI6eyJuYW1lIjoiZGFzaGVkIn19fV0sWzAsMywiMV9CIiwxXSxbMCwyLCJnIiwyXV0=
%	\[\begin{tikzcd}
%		& B & \\
%		& A \\
%		A && B
%		\arrow["g"', dashed, from=1-2, to=2-2]
%		\arrow["g"', from=1-2, to=3-1]
%		\arrow["{1_B}"{description}, from=1-2, to=3-3]
%		\arrow["{1_A}", from=2-2, to=3-1]
%		\arrow["{\pi_2}"', from=2-2, to=3-3]
%	\end{tikzcd}\]
%	
%	and so $\pi_2 = 1_B$. 
	
	

%	
%	We will now show that $B$ is terminal. We must show that there is a unique map $A \to B$ and $B \to B$. 
%	
%	We know $\pi_2$ is unique. So, we must show that there are no more maps $B \to B$. Suppose there is some $g: B \to B$. If $B \times B = A$, $A$ being initial forces any $f: B \to A$ to be a right inverse for $\pi_2$.
%	
%	% https://q.uiver.app/#q=WzAsNCxbMSwwLCJCIl0sWzEsMiwiQSJdLFswLDMsIkEiXSxbMiwzLCJCIl0sWzEsMiwiMV9BIl0sWzEsMywiXFxwaV8yIiwyXSxbMCwxLCI/IiwwLHsic3R5bGUiOnsiYm9keSI6eyJuYW1lIjoiZGFzaGVkIn19fV0sWzAsMywiMV9CIl0sWzAsMiwiPyIsMix7InN0eWxlIjp7ImJvZHkiOnsibmFtZSI6ImRhc2hlZCJ9fX1dXQ==
%	\[\begin{tikzcd}
%		& B & \\
%		\\
%		& A \\
%		A && B
%		\arrow["{?}", dashed, from=1-2, to=3-2]
%		\arrow["{?}"', dashed, from=1-2, to=4-1]
%		\arrow["{1_B}", from=1-2, to=4-3]
%		\arrow["{1_A}", from=3-2, to=4-1]
%		\arrow["{\pi_2}"', from=3-2, to=4-3]
%	\end{tikzcd}\]
%	
%	but similarly,  $g$ must be $1_B$:
%	% https://q.uiver.app/#q=WzAsNCxbMSwwLCJCIl0sWzEsMiwiQSJdLFswLDMsIkEiXSxbMiwzLCJCIl0sWzEsMiwiMV9BIl0sWzEsMywiXFxwaV8yIiwyXSxbMCwxLCI/IiwwLHsic3R5bGUiOnsiYm9keSI6eyJuYW1lIjoiZGFzaGVkIn19fV0sWzAsMywiZyJdLFswLDIsIj8iLDIseyJzdHlsZSI6eyJib2R5Ijp7Im5hbWUiOiJkYXNoZWQifX19XV0=
%	\[\begin{tikzcd}
%		& B & \\
%		\\
%		& A \\
%		A && B
%		\arrow["{?}", dashed, from=1-2, to=3-2]
%		\arrow["{?}"', dashed, from=1-2, to=4-1]
%		\arrow["g", from=1-2, to=4-3]
%		\arrow["{1_A}", from=3-2, to=4-1]
%		\arrow["{\pi_2}"', from=3-2, to=4-3]
%	\end{tikzcd}\]
%	
%	If instead, $B \times B = B$, the same logic as we've used above shows that $p_2: B \to B$ is $1_B$ (think about the unique morphism induced by the maps $p_1, 1_B$). Also similar to above, $p_1$ is unique. Again similar to above, we conclude $1_B$ is the 

	

%%	Consider the category with two elements and an isomorphism between them. Then, neither obj
%
%$p_A \circ f = 1_A. f \circ p_A = p_A \circ g$. 
%
%$g \circ f \circ p_A = p_A \circ f \circ g$
%
%$g \circ p_A \circ g = g$
%
%
%
%
%$f \circ p_A = p_A \circ g$
%$g \circ p_A = p_A \circ f$
%
%$g \circ f \circ p_A = g \circ p_A \circ g = p_A \circ f \circ g = g$
%
%$f \circ p_A \circ g = p_A \circ g \circ g$



\end{solution}





\begin{problem}[93]
	Let $\Small$ be the category with one object $*$ and morphisms $\Small(*, *) = \{1, f\}$ where $f \circ f = f$. Let $D : \Small \to \Set$ be a diagram. Use the method of the "products and equalizers give all limits" proof to compute the limit of $D$ explicitly.
\end{problem}
\begin{solution}
	
	The aforementioned proof built the product over all objects in $\Small$, then the product over all morphisms in $\Small$, found two morphisms between them (one that bypassed morphisms $f: S \to T$, and one that respected them) then equalized them. The limit cone was then $L \xrightarrow{\lambda} \prod_S DS \xrightarrow{\pi_S} DS$. 
	
	We have one object, so the product over all objects in $\Small$ is $A := D*$ with projection $1_*$ (if the projection were $f$, there would be no unique morphism associated to the map $1_*$). 
	
	The product over all $g: * \to *$ looks like $A \times A$ with projections $\pi_1$ and $\pi_f$. 
	
	Write $g$ to be either $1$ or $f$. 
	
	% https://q.uiver.app/#q=WzAsMyxbMCwwLCJBIl0sWzIsMCwiQVxcdGltZXMgQSJdLFsyLDIsIkEiXSxbMSwyLCJcXHBpXzEgKFxcdGV4dHtXTE9HfSkiXSxbMCwyLCIxX0EiLDJdLFswLDEsIlxcYWxwaGEiXV0=
	\[\begin{tikzcd}
		A && {A\times A} \\
		\\
		&& A
		\arrow["\alpha", from=1-1, to=1-3]
		\arrow["{1_A}"', from=1-1, to=3-3]
		\arrow["{\pi_g}", from=1-3, to=3-3]
	\end{tikzcd}\]
	
	tells us that $\pi_g \alpha = 1_A$. That is, $\alpha = (1, 1)$.
	
	% https://q.uiver.app/#q=WzAsNCxbMCwwLCJBIl0sWzAsMiwiQSJdLFsyLDIsIkEiXSxbMiwwLCJBIFxcdGltZXMgQSJdLFswLDEsIjFfQSJdLFsxLDIsImciXSxbMCwzLCJcXGJldGEiXSxbMywyLCJcXHBpXzEiLDJdXQ==
	\[\begin{tikzcd}
		A && {A \times A} \\
		\\
		A && A
		\arrow["\beta", from=1-1, to=1-3]
		\arrow["{1_A}", from=1-1, to=3-1]
		\arrow["{\pi_g}", from=1-3, to=3-3]
		\arrow["Dg", from=3-1, to=3-3]
	\end{tikzcd}\]
	
	tells us that $\pi_g \beta = Dg$. That is, $\beta = (1, Df)$
		
	The equalizer in $\Set$ is $\{x : \alpha(x) = \beta(x)\}$ equipped with the inclusion map $\iota$. In this case, we get $L = \{x : (x, x) = (x, Df(x))\}$. Equality in the first  slot is always true, so $L = \{x \in A : x = Df(x)\}$ is the set of fixed points of $Df$. The proof that products and equalizers give all limits proves that $(L, \iota)$ is a limit cone. Hence, $L$ is the limit of $D$.  \qed
	
	
\end{solution}





\begin{problem}[96, 97]
	Prove $\Delta\_ \dashv \lim$ (I want to do this proof in full so that I understand it)
\end{problem}
\begin{solution}
	
	NB: $\lim: [\Small, \mc{C}] \to \mc{C}, \Delta_\_: \mc{C} \to [\Small, \mc{C}]$
	
	We begin by defining $\eta: 1_\mc{C} \natt \lim \circ \Delta_\_$ as our unit such that $\eta_A : A \to \lim \Delta_A$ uniquely induced by choosing the morphisms $1_A$, as in 
	% https://q.uiver.app/#q=WzAsMyxbMCwwLCJBIl0sWzAsMSwiXFxsaW0gXFxEZWx0YV9BIl0sWzIsMSwiQSA9IFxcRGVsdGFfQShTKSJdLFswLDEsIlxcZXhpc3RzISIsMix7InN0eWxlIjp7ImJvZHkiOnsibmFtZSI6ImRhc2hlZCJ9fX1dLFsxLDIsInBfUyIsMl0sWzAsMiwiMV9BIl1d
	\[\begin{tikzcd}
		A && \\
		{\lim \Delta_A} && {A = \Delta_A(S)}
		\arrow["{\exists!}"', dashed, from=1-1, to=2-1]
		\arrow["{1_A}", from=1-1, to=2-3]
		\arrow["{p_S}"', from=2-1, to=2-3]
	\end{tikzcd}\]
	
	(where there are as many $S$'s as objects in $\Small$ and so we trivially have a cone). That is, $p_S \circ \eta_A = 1_A$ for every projection $p_S$. 
	
	For $\eta$ to be natural, the following must commute for each $f: A \to B$ in $\mc{C}$. 
	% https://q.uiver.app/#q=WzAsNCxbMCwwLCJBIl0sWzEsMCwiQiJdLFsxLDEsIlxcbGltIFxcRGVsdGFfQiJdLFswLDEsIlxcbGltIFxcRGVsdGFfQSJdLFswLDMsIlxcZXRhX0EiLDJdLFswLDEsImYiXSxbMSwyLCJcXGV0YV9CIl0sWzMsMiwiXFxsaW0gXFxEZWx0YV9mIiwyXV0=
	\[\begin{tikzcd}
		A & B \\
		{\lim \Delta_A} & {\lim \Delta_B}
		\arrow["f", from=1-1, to=1-2]
		\arrow["{\eta_A}"', from=1-1, to=2-1]
		\arrow["{\eta_B}", from=1-2, to=2-2]
		\arrow["{\lim \Delta_f}"', from=2-1, to=2-2]
	\end{tikzcd}\]
	
	Let the $B$ limit have projections $\pi_S$. Then, $\pi_S \circ \eta_B \circ f = f$, and $\pi_S \circ \lim \Delta_f \circ \eta_A = (\Delta_f)_S \circ p_S \circ \eta_A$ by the commutative square from the lemma. This  $= (\Delta_f)_S = f$. Then, because the projections all agree, by uniqueness of morphisms, $\eta_B \circ f = \lim \Delta_f \circ \eta_A$ and so $\eta$ is natural!
	
	We now define $\eps: \Delta\_ \circ \lim \natt 1_{[\Small, \mc{C}]}$ with components $\eps_F : \Delta_{\lim F} \to F$. $\Delta_{\lim F}$ itself has components given by $\lim F$ for each component. We define \[\eps_F((\Delta_{\lim F})_S) = \lambda_S: \lim F \to FS\] the projections from the definition of the limit. 
	
	Then, for $\eps$ to be natural, the following must commute:
	
	% https://q.uiver.app/#q=WzAsOCxbMCwwLCJcXGxpbSBGIl0sWzMsMCwiXFxsaW0gRyJdLFswLDMsIkYiXSxbMywzLCJHIl0sWzEsMSwiXFxsaW0gRiJdLFsxLDIsIkZTIl0sWzIsMiwiR1MiXSxbMiwxLCJcXGxpbSBHIl0sWzAsMiwiXFxlcHNfRiIsMix7ImxldmVsIjoyfV0sWzIsMywiXFxhbHBoYSIsMix7ImxldmVsIjoyfV0sWzAsMSwiXFxsaW0gXFxhbHBoYSIsMCx7ImxldmVsIjoyfV0sWzEsMywiXFxlcHNfRyIsMCx7ImxldmVsIjoyfV0sWzQsNSwicF9TIiwwLHsic3R5bGUiOnsidGFpbCI6eyJuYW1lIjoibWFwcyB0byJ9fX1dLFs1LDYsIlxcYWxwaGFfUyIsMCx7InN0eWxlIjp7InRhaWwiOnsibmFtZSI6Im1hcHMgdG8ifX19XSxbNCw3LCJcXGxpbSBcXGFscGhhIiwwLHsic3R5bGUiOnsidGFpbCI6eyJuYW1lIjoibWFwcyB0byJ9fX1dLFs3LDYsIlxccGlfUyIsMix7InN0eWxlIjp7InRhaWwiOnsibmFtZSI6Im1hcHMgdG8ifX19XV0=
	\[\begin{tikzcd}
		{\lim F} &&& {\lim G} \\
		& {\lim F} & {\lim G} \\
		& FS & GS \\
		F &&& G
		\arrow["{\lim \alpha}", Rightarrow, from=1-1, to=1-4, nfold]
		\arrow["{\eps_F}"', Rightarrow, from=1-1, to=4-1, nfold]
		\arrow["{\eps_G}", Rightarrow, from=1-4, to=4-4, nfold]
		\arrow["{\lim \alpha}", maps to, from=2-2, to=2-3]
		\arrow["{\lambda_S}", maps to, from=2-2, to=3-2]
		\arrow["{\pi_S}"', maps to, from=2-3, to=3-3]
		\arrow["{\alpha_S}", maps to, from=3-2, to=3-3]
		\arrow["\alpha"', Rightarrow, from=4-1, to=4-4, nfold]
	\end{tikzcd}\]
	
	However, the lemma tells us that tracing any component of $\lim F$ around the diagram works! Hence, $\eps$ is natural. 
	
	Now, we prove the triangle identities. 
	
	We know that $\lambda_S \circ \lim(\eps_F) \circ \eta_{\lim F} = (\eps_F)_S \circ p_S \circ \eta_{\lim F} = \lambda_S \circ p_S \circ \eta_{\lim F} = \lambda_S \circ 1_{\lim F} = \lambda_S$. 
	
	So, the outer part of the following diagram commutes. By uniqueness and the fact that $1$ works, we get that the upper triangle commutes. The upper triangle is the $U$ triangle identity
% https://q.uiver.app/#q=WzAsNCxbMCwwLCJcXGxpbSBGIl0sWzQsMCwiXFxsaW0gXFxEZWx0YV97XFxsaW0gRn0iXSxbNCw0LCJcXGxpbSBGIl0sWzQsNSwiRlMiXSxbMCwxLCJcXGV0YV97XFxsaW0gRn0iXSxbMSwyLCJcXGxpbSAoXFxlcHNfRikiXSxbMiwzLCJcXGxhbWJkYV9TIl0sWzAsMywiXFxsYW1iZGFfUyIsMix7ImN1cnZlIjo1fV0sWzAsMiwiMSIsMSx7InN0eWxlIjp7ImJvZHkiOnsibmFtZSI6ImRhc2hlZCJ9fX1dXQ==
\[\begin{tikzcd}
	{\lim F} &&&& {\lim \Delta_{\lim F}} \\
	\\
	\\
	\\
	&&&& {\lim F} \\
	&&&& FS
	\arrow["{\eta_{\lim F}}", from=1-1, to=1-5]
	\arrow["1"{description}, dashed, from=1-1, to=5-5]
	\arrow["{\lambda_S}"', curve={height=30pt}, from=1-1, to=6-5]
	\arrow["{\lim (\eps_F)}", from=1-5, to=5-5]
	\arrow["{\lambda_S}", from=5-5, to=6-5]
\end{tikzcd}\]



The other triangle identity gives
% https://q.uiver.app/#q=WzAsMyxbMCwwLCJcXERlbHRhX0YiXSxbMiwyLCJcXERlbHRhX0YiXSxbMiwwLCJcXERlbHRhX3tcXGxpbSBcXERlbHRhX0Z9Il0sWzAsMSwiMSIsMix7ImxldmVsIjoyfV0sWzAsMiwiXFxEZWx0YV97XFxldGFfRn0iLDAseyJsZXZlbCI6Mn1dLFsyLDEsIlxcZXBzX3tcXERlbHRhX0Z9IiwwLHsibGV2ZWwiOjJ9XV0=
\[\begin{tikzcd}
	{\Delta_F} && {\Delta_{\lim \Delta_F}} \\
	\\
	&& {\Delta_F}
	\arrow["{\Delta_{\eta_F}}", Rightarrow, from=1-1, to=1-3, nfold]
	\arrow["1"', Rightarrow, from=1-1, to=3-3, nfold]
	\arrow["{\eps_{\Delta_F}}", Rightarrow, from=1-3, to=3-3, nfold]
\end{tikzcd}\]
Taking components, we get
% https://q.uiver.app/#q=WzAsMyxbMCwwLCJGIl0sWzIsMCwiXFxsaW0gXFxEZWx0YV9GIl0sWzIsMiwiRiJdLFswLDEsIlxcZXRhX0YiXSxbMSwyLCJwX1MiXSxbMCwyLCIxIiwyXV0=
\[\begin{tikzcd}
	F && {\lim \Delta_F} \\
	\\
	&& F
	\arrow["{\eta_F}", from=1-1, to=1-3]
	\arrow["1"', from=1-1, to=3-3]
	\arrow["{p_S}", from=1-3, to=3-3]
\end{tikzcd}\]

because $\eps_{\Delta_F}: \Delta_{\lim \Delta_F} \to \Delta_F$ has $(\eps_{\Delta_F})_S = p_S : \lim \Delta_F \to \Delta_F(S) = F$, which commutes by our construction of $p_S$ from our construction of $\eta$. So, the other triangle identity is satisfied. 

Thus, we have an adjunction! \qed






%	
	
	
%	Hence, it also comes equipped with $\lambda_S: \lim F \to FS$ for each $S$. 
	
	
	
	
	
	
	
\end{solution}

\end{document}
